% ============================================================================
%  Neutrino fog for all materials and comparison with CRESST
%  Figures: fog_all_materials.png, fog_vs_cresst.png (NuFog_v2/outputs)
%  New bibliography entries at the end of the file.
% ============================================================================

\section{Neutrino fog across target materials and comparison with CRESST}
\label{sec:fog_cresst}

\subsection{Discovery limits for all target materials}

Figure~\ref{fig:fog_all} collects the $3\sigma$ discovery limits for an exposure of
$1$\,t$\cdot$yr for the five crystals studied in this work, at energy thresholds of
$30$\,meV and $100$\,meV (with resolution $\sigma = E_{\rm th}/5$) and $1$\,eV. The
curves define the cross section below which a DM signal can no longer be distinguished
from coherent elastic neutrino--nucleus scattering, given the flux uncertainties of
Sec.~\ref{sec:likelihood}.

Two regimes are visible.
\begin{itemize}
  \item \textbf{Above $m_\chi \simeq 0.3$\,GeV} the curves of all materials nearly
  coincide. Here both the DM and the neutrino signals are dominated by nuclear recoils
  in the impulse regime. The DM rate and the CEvNS rate then scale with the same
  coherent enhancement, so the target composition largely cancels. The feature near
  $0.5$\,GeV is where the DM recoil spectrum best mimics the $^8$B solar-neutrino
  spectrum.
  \item \textbf{Below $m_\chi \simeq 0.1$\,GeV} the limits are set by the lattice
  response and separate by up to an order of magnitude. At $E_{\rm th}=30$\,meV,
  Al$_2$O$_3$ and Si give the lowest discovery limits. The heavier targets TeO$_2$ and
  CaWO$_4$ are weaker by factors of $8$ and $16$, because their phonon spectra put less
  weight at high frequency and the momentum transfer is shared with heavy nuclei. At
  $E_{\rm th}=1$\,eV the low-mass reach ends at $m_\chi \simeq 20$--$40$\,MeV for all
  materials, and the shoulder of the CaWO$_4$ and TeO$_2$ curves near $0.1$\,GeV marks
  the transition between their light (O) and heavy (W, Te) constituents.
\end{itemize}

\begin{figure}[t]
  \centering
  \includegraphics[width=\textwidth]{fog_all_materials.png}
  \caption{$3\sigma$ discovery limits on the spin-independent DM--nucleon cross section
  in the presence of the neutrino background, for an exposure of 1\,t$\cdot$yr, for
  CaWO$_4$, Al$_2$O$_3$, TeO$_2$, LiNbO$_3$ and Si. Left: $E_{\rm th}=30$\,meV;
  centre: $E_{\rm th}=100$\,meV, both with Gaussian resolution $\sigma=E_{\rm th}/5$;
  right: $E_{\rm th}=1$\,eV.}
  \label{fig:fog_all}
\end{figure}

\subsection{Comparison with CRESST}

Figure~\ref{fig:fog_cresst}a places these results in the context of the CRESST
experiment. The current CRESST-III limits from the CaWO$_4$ ``Detector A''
(threshold $30.1$\,eV)~\cite{CRESST2019} and the silicon-on-sapphire (SOS) detector
(threshold $6.7$\,eV)~\cite{CRESST2025Perspective}, together with the projected
sensitivity of the CRESST upgrade~\cite{CRESST2025Perspective}, are compared with the
neutrino fog computed here for a $30$\,meV threshold. The upgrade projection assumes a
100-fold reduction of the low-energy excess and $515$\,kg$\cdot$day of exposure, split
between 2\,g modules with a $5$\,eV threshold and 24\,g modules with a $30.1$\,eV
threshold.

Several conclusions follow.
\begin{itemize}
  \item At $m_\chi \simeq 1$\,GeV the current CRESST-III limit lies about six orders of
  magnitude above the neutrino fog, and the upgrade projection about two. This matches
  the conclusion of Ref.~\cite{CRESST2025Perspective}: CRESST will remain
  background-free with respect to neutrinos for the foreseeable future, and the
  low-energy excess, not the neutrino background, is what limits its sensitivity.
  \item The lowest mass probed by CRESST, $73$\,MeV with the SOS detector, is set by a
  threshold of several eV. Thresholds of $\mathcal{O}(10$--$100)$\,meV, where
  multiphonon emission dominates, extend the region in which the neutrino fog becomes
  the relevant background down to $m_\chi \simeq 0.5$\,MeV. This region is not yet
  accessible to any CRESST-type detector.
  \item In this low-mass region the choice of target matters most
  (Fig.~\ref{fig:fog_all}). CaWO$_4$, the CRESST baseline material, gives the weakest
  discovery limit below $50$\,MeV, while Al$_2$O$_3$ and Si, which CRESST already
  operates, gain more than an order of magnitude.
\end{itemize}

Figure~\ref{fig:fog_cresst}b validates the calculation in the regime studied by CRESST.
For CaWO$_4$ with $E_{\rm th} = 1$\,eV and an exposure of $1$\,t$\cdot$yr, our discovery
limit reproduces the shape of the corresponding curve of Bento \emph{et
al.}~\cite{Bento2024}, including the $^8$B feature near $0.5$\,GeV and the shoulder near
$0.1$\,GeV. Our limit is lower by a factor of $0.61$--$0.78$ over
$0.05$--$4$\,GeV. Part of this difference comes from the digitization of the reference
curve, which reproduces the published CRESST-III 2019 limit to within $10$--$20\%$. The
rest is consistent with differences in the neutrino-flux inputs and in the treatment of
the low-energy response. The other materials at the same threshold show the same
convergence at high mass and the same separation at low mass discussed above.

\begin{figure}[t]
  \centering
  \includegraphics[width=\textwidth]{fog_vs_cresst.png}
  \caption{(a) CRESST-III exclusion limits at 90\% CL from Detector A
  (2019)~\cite{CRESST2019} and from the SOS detector (2024), and the projected
  sensitivity of the CRESST upgrade with 100 times less low-energy excess (median, envelope of the 5\,eV and
  30.1\,eV module sets)~\cite{CRESST2025Perspective}, compared with the neutrino fog of
  this work for $E_{\rm th}=30$\,meV, $\sigma = E_{\rm th}/5$ and 1\,t$\cdot$yr. The grey
  region is the CaWO$_4$ neutrino fog of Ref.~\cite{Bento2024}.
  (b) Discovery limits for $E_{\rm th}=1$\,eV and 1\,t$\cdot$yr: this work for all
  materials (solid) and Bento \emph{et al.}\ for CaWO$_4$ (dotted). Curves other than the
  CRESST-III 2019 limit were digitized from Ref.~\cite{CRESST2025Perspective}.}
  \label{fig:fog_cresst}
\end{figure}

% ---------------- BibTeX entries (move to your .bib) ------------------------
% @article{CRESST2019,
%   author  = {Abdelhameed, A. H. and others},
%   collaboration = {CRESST},
%   title   = {First results from the {CRESST-III} low-mass dark matter program},
%   journal = {Phys. Rev. D}, volume = {100}, pages = {102002}, year = {2019},
%   eprint  = {1904.00498}, archivePrefix = {arXiv}}
%
% @article{CRESST2025Perspective,
%   author  = {Angloher, G. and others},
%   collaboration = {CRESST},
%   title   = {The {CRESST} experiment towards the next generation of sub {GeV} direct dark matter detection},
%   journal = {Commun. Phys.}, year = {2025},
%   doi     = {10.1038/s42005-025-02476-5}}
%
% @article{Bento2024,
%   author  = {Bento, A. and others},
%   title   = {Solar neutrinos in cryogenic detectors},
%   journal = {Eur. Phys. J. C}, volume = {84}, pages = {1118}, year = {2024}}
